\part{Alternative Layouts}

\chapter{Dvorak}
\label{ch:dvorak}

The Dvorak Simplified Keyboard was patented in the 1930s by August Dvorak and William Dealey. \needs{patent year and details} It is the best-known alternative to QWERTY and the layout against which the curricula of Part~VII are compared.

\section{The Layout}

In the standard arrangement the rows are as follows, with punctuation shown as it appears.

\begin{center}
\begin{tabular}{@{}lll@{}}
\toprule
Row & QWERTY & Dvorak \\
\midrule
Top & \texttt{q w e r t y u i o p} & \texttt{' , . p y f g c r l} \\
Home & \texttt{a s d f g h j k l ;} & \texttt{a o e u i d h t n s} \\
Bottom & \texttt{z x c v b n m , . /} & \texttt{; q j k x b m w v z} \\
\bottomrule
\end{tabular}
\end{center}

Dvorak places all five vowels and the frequent consonants D, H, T, N, S on the home row, and the design principles usually cited are alternation between hands and concentration of work on the strongest fingers. Only two letters, A and M, occupy the same position in both layouts.

\section{The Disputed Record}

Claims that Dvorak is markedly faster have been repeated for decades and contested with equal persistence, and the evidence includes mid-century studies and later reappraisals whose conclusions differ. \needs{cite the principal studies and reappraisals without endorsing either side} This book does not adjudicate. The comparison of Part~VII is designed so that it does not depend on the outcome.

\section{Position Matching}

Because both layouts are laid out on the same physical keys, any QWERTY lesson can be transposed to Dvorak by position: the key in each physical location is replaced by the key that occupies that location in Dvorak. The letters of the mapping are these.

\begin{center}
\begin{tabular}{@{}lll@{}}
\toprule
Row & QWERTY key & Dvorak key at that position \\
\midrule
Top & Q W E R T Y U I O P & \texttt{'} \texttt{,} \texttt{.} P Y F G C R L \\
Home & A S D F G H J K L ; & A O E U I D H T N S \\
Bottom & Z X C V B N M , . / & \texttt{;} Q J K X B M W V Z \\
\bottomrule
\end{tabular}
\end{center}

Applied to the Academy Exercises 1, 2, 6 and 13 of Chapter~8, the mapping gives the following.

\begin{center}
\begin{tabular}{@{}clll@{}}
\toprule
Exercise & QWERTY new keys & Dvorak by position & Consequence \\
\midrule
1 & A S D F ; L K J & A O E U S N T H & vowels and four consonants; no semicolon \\
2 & G H & I D & completes the home row \\
6 & W I & comma, C & one new key is punctuation \\
13 & Z & semicolon & a lesson with no letter at all \\
\bottomrule
\end{tabular}
\end{center}

This is the cost of position matching, and it is the subject of Chapter~24: the same geometry yields a different lexicon at each stage, and at some stages no lexicon at all.

\section{Matching Letters Instead}

The alternative is to match content, teaching the same letters in the same order and letting the geometry differ. That preserves the vocabulary $\Vocab_i$ at each stage but breaks the mapping between lesson number and keyboard region. Chapters 24 and 25 present both and compare them.

\chapter{Chorded Input and the Unit of Action}

Chapter~3 described where machine shorthand came from. This chapter asks a different question: what happens to a tutor when the primitive action is no longer a single key?

\section{Chords and Strokes}

In a chorded system a character or syllable is produced by pressing several keys at once. The unit of action is the \emph{stroke}, a set of keys pressed and released together. Let $k$ be the number of keys and $s\subseteq\{1,\dots,k\}$ a stroke. The number of distinct strokes available on a board of $k$ keys is $2^k-1$, which is why a small chorded keyboard can address a large vocabulary.

\section{What Changes for the Curriculum}

In a single-key tutor the lesson order is an order over keys. In a chorded tutor it is an order over strokes, and strokes have structure: they overlap, so learning one stroke changes the cost of learning another. A curriculum can exploit the structure, for example by introducing strokes that share fingers together.

The error model changes as well. A wrong key in a single-key system is one wrong character. A wrong chord may differ from the intended one by a single finger, so the natural error metric is a distance between sets, for instance
\[
d(s,s')=|s\,\triangle\,s'|,
\]
the number of keys by which two strokes differ. The learner state of Chapter~17 has to be redefined with strokes in place of keys, and the transition term becomes a matrix over pairs of strokes.

\section{Phonetic Briefs}

Court-reporting theories also compress high-frequency words and phrases into single strokes, called briefs. A tutor for such a system teaches a dictionary as well as a motor skill. \needs{a specific stenotype theory and its learning sequence}

\section{The Pedagogical Primitive}

The finding of this chapter is a point of method. Single-key tutors treat the key as the unit of both action and error. The unit of action can differ from the unit of the tutor's display, so the definition of the tutoring primitive is a design decision that every tutor makes, whether or not it says so.

\chapter{One-Handed and Ergonomic Layouts}

The relation between hand and keyboard is not fixed, and a substantial body of designs reshapes it.

\section{Split and Staggered Boards}

The row stagger of a standard keyboard descends from the mechanics of the typewriter's linkages, not from the geometry of the hand. Split keyboards separate the halves, and columnar layouts align keys with the movement of the fingers. \needs{representative designs and dates}

\section{One-Handed Arrangements}

One-handed layouts serve users with a single functioning hand and can also be used to reduce load. Some arrangements fold a standard layout onto one half, so that the other hand's keys are reached by a modifier. \needs{sources on one-handed layouts, including designs derived from QWERTY}

\section{Implications for Tutors}

Each layout implies a curriculum, and the curriculum must be regenerated when the geometry changes. This is the argument for treating the tutor's ladder as a template with three parameters: the key-introduction order, the lexicon, and the layout header, so that a new layout is a new instantiation and not a new tutor.

\chapter*{Study Questions, Part IV}
\addcontentsline{toc}{chapter}{Study Questions, Part IV}
\begin{enumerate}
\item Transpose Exercise~2 of the Academy pages to Dvorak by position, and state the new keys.
\item Why does position matching produce an empty $\Vocab_i^{\mathrm{new}}$ at some exercises? Give the two exercises in the design order.
\item On a chorded board with $k=6$ keys, how many strokes exist? What is $d(s,s')$ for strokes $\{1,2,4\}$ and $\{1,3,4\}$?
\item Why can a tutor for a chorded system not treat the key as its unit of error?
\end{enumerate}
